\answer{ex:03:1}{(а)~$35\to17.5\mV$ при $g_E=g_L$; $56\to40\mV$ при $g_E=4g_L$ (вычитание дало бы $38.5\mV$); шунт равносилен делению $g_E$ на $3$. (б)~$\tau_{\text{эфф}}=10\to5\ms$, окно $6.9\to3.5\ms$. (в)~$v^*=2$~мм/с.}
\answer{ex:03:2}{(а)~$\operatorname{tr}=\frac{w_{EE}-1}{\tau_E}-\frac1{\tau_I}$, $\det=\frac{1-w_{EE}+w_{EI}w_{IE}}{\tau_E\tau_I}$. (б)~$\omega=\sqrt{\det}$; граница $\tau_I=20\ms$, период $73\ms$; $\lambda=(-0.225\pm0.156i)\,\text{мс}^{-1}$ при $\tau_I=2\ms$ и $(0.0083\pm0.070i)\,\text{мс}^{-1}$ при $30\ms$ (период $89\ms$), рост ограничивает $[\cdot]_+$. (в)~$-1/3$; при $w_{EE}>1$.}
\answer{ex:03:3}{(б)~$\omega=\sqrt{G^2-1}/\tau_E$, $d_c=\tau_E\arccos(-1/G)/\sqrt{G^2-1}$, период $2\pi\tau_E/\sqrt{G^2-1}\in(2d_c,4d_c)$. (в)~$G=2$: $d_c=12.1\ms$, период $36\ms$; $G=5$: $d_c=3.6\ms$, период $12.8\ms$.}
\answer{ex:03:4}{(а)~$0.73$ и $0.53$; $1$ и $0$; после перестановки $0.9$ и $0$. (б)~Переключения при $I_2=\beta I_1=2$ и $I_2=I_1/\beta=0.5$, ширина петли $1.5$. (в)~$+\tau\ln10/(\beta-1)=2.3\tau$ на каждый десятичный порядок $\varepsilon$.}
\answer{ex:03:5}{(а)~Три посредника, $d_k=5$, $10$, $15\ms$; $k\,T_I\ge15\ms$. (в)~Вне полосы детектор отвечает в обоих направлениях; при $T_I=10\ms$ хватает двух.}
\answer{ex:03:6}{(а)~$n+M+1$ нейронов ($M$ --- число наборов с $f=1$); для $x\oplus y\oplus z$ --- $8$. (б)~\nt{ГАМК} $=[x+y+z\ge2]$, выход $=[x+y+z-2\,\nt{ГАМК}\ge1]$. (в)~Прямое возбуждение надо задержать на время пути через \nt{ГАМК}; при неизвестном разбросе моментов --- нельзя.}
\answer{ex:03:7}{(а)~Один посредник, возбуждаемый всеми и тормозящий всех, плюс самовозбуждение $P=\beta I$. (б)~$\binom84=70<100\le\binom94=126$: $9$ посредников. (в)~$100$: ранг $\beta(J-I)$ равен $n$.}
\answer{ex:03:8}{(а)~$h_c=\frac1{2w_{EE}}+\frac{w_{EI}w_{IE}-w_{EE}}{4w_{EE}^2}=\frac7{18}\approx0.39$; условие~(ii) нарушается при $h=4$. (б)~Смена знака между $h=0.38$ и $0.40$; при резком включении $h\gtrsim2.4$ из тишины активность уходит в разнос.}
